\subsection{Two uniform Frobenius tails by one alternant}
\label{sec:flag-tail-alternants}
We supply the two numerical remainder estimates in the tower comparison
without the polynomial recursions of \texttt{RestDegreePowerBound} and
\texttt{TestPolynomialRemainderBound}. All source comparisons in this
subsection refer to the pinned revision \texttt{ecd09ff9} of
\texttt{jsmorph/rh}. The estimates concern the original remainder, not
a substitute geometric model.

Put \(N=p+q,\ 3q+1\le N,\ q\ge1,\ f=|\Psi|\ge2\),
\(d=fN(N-1)/2\), and \(d_0=fpq\). Use the original degree functional
\(\deg_{\rm Fl}=(-1)^d\operatorname{alt}_{\rho}\), where
\(\rho_{\tau i}=N-1-i\) and \(\operatorname{alt}_a\) is signed
coefficient extraction at all permutations of the exponent lists \(a\).
Put \(p_{\tau k}=\sum_i X_{\tau i}^k\).
Every homogeneous polynomial \(F\) of degree \(d+1\) has a decomposition
\[
 F=\sum_{\tau}\sum_{k=1}^{N}p_{\tau k}Q_{\tau k},
 \qquad \deg Q_{\tau k}=d+1-k.
 \tag{FA1}
\]
Indeed, modulo the power-sum ideal the elementary symmetric functions
of each label vanish by Newton's identities. In one label, numbered
\(1,\ldots,N\), the complete symmetric function
\(h_i(X_i,\ldots,X_N)\) then vanishes: its generating series reduces
to \(\prod_{j<i}(1-X_jt)\). Its lexicographic leading monomial is
\(X_i^i\). Reducing by these monic homogeneous polynomials, in each
label, leaves monomials of total degree at most \(d\).
Homogeneity now gives (FA1).

\paragraph{The extraction identity.}
Define \(\lambda_{\tau k}=\sum_{i=0}^{N-1}
\operatorname{alt}_{\rho+k e_{\tau i}}\). For \(1\le k,m\le N\),
\[
 \lambda_{\tau k}(p_{\sigma m}G)
 ={\bf1}_{\tau=\sigma,\ k=m}\, k\operatorname{alt}_{\rho}(G).
 \tag{FA2}
\]
Here and below a coefficient of negative exponent is zero.
For a direct proof use the exterior product of the exponent basis
\(e_{N-1}\wedge\cdots\wedge e_0\) in each label. Multiplication of
the polynomial by a power sum acts on coefficient extraction
by the derivation induced by the downward shift \(L^m\), where
\(Le_0=0,\ Le_a=e_{a-1}\) for \(a>0\).
Bumping a single exponent and summing acts by the derivation of
\(S^k\), where \(Se_a=e_{a+1}\).
The downward derivation kills the original exterior product.
Different labels commute. In one label the commutator
\([L^m,S^k]e_a\) is zero except for \(a<m\le a+k\), when it equals
\(e_{a+k-m}\). For \(k\ne m\) this is a different exponent
between \(0\) and \(k-1\le N-1\), already present in the wedge.
For \(k=m\) there are \(k\) identical contributions to the
original wedge. This proves (FA2), and the same downward-shift
argument gives
\(\operatorname{alt}_{\rho}(p_{\sigma m}G)=0\) for every \(m>0\).
Consequently, for \(A\) homogeneous of degree \(k-1\),
\[
 k\deg_{\rm Fl}(A Q_{\tau k})
 =(-1)^d\lambda_{\tau k}(A F).
 \tag{FA3}
\]
In particular these numbers are independent of the choice in (FA1).

For a vector \(z=(z_{\tau i})\), define
\[
 P(z)=\operatorname{alt}_{\rho}\exp\Big(\sum_{\tau i}z_{\tau i}X_{\tau i}\Big)
     =\prod_\tau\det\left(\frac{z_{\tau i}^{N-1-j}}{(N-1-j)!}\right)_{j,i}.
\]
It is the full Vandermonde times \(\prod_{j=0}^{N-1}(j!)^{-f}\).
The bumped determinant yields
\[
 \lambda_{\tau k}\exp\Big(\sum_{\sigma i}z_{\sigma i}X_{\sigma i}\Big)
 =P(z)H_{\tau k}(z),\qquad
 H_{\tau k}(z)=
 \sum_{r=0}^{k-1}(-1)^r\frac{(N-1-r)!}{(N-1-r+k)!}
 s_{(k-r,1^r)}(z_\tau).
 \tag{FA4}
\]
Indeed a bumped exponent which stays below \(N\) gives a repeated
row. For the remaining rows the bump is moved to the first row
through \(r\) rows; the ratio of alternants is the indicated hook
Schur polynomial. The factorial ratio is retained.
The hook-content formula (or the Vandermonde limit in the
Schur alternant) gives
\[
 \frac{(N-1-r)!}{(N-1-r+k)!}s_{(k-r,1^r)}(1^N)
   =\frac{\binom{k-1}{r}}{k!}.
\]
Since Schur polynomials have nonnegative monomial coefficients
(the semistandard-tableau expansion),
\(\|H_{\tau k}\|_{\ell^1}\le 2^{k-1}/k!\).

\paragraph{Regular-weight remainder.}
Let \(u\) be strictly dominant, \(U=\max_{\tau i}|u_{\tau i}|\),
\(W_\tau=\max_i u_{\tau i}-\min_i u_{\tau i}\),
\(W=\sum_\tau W_\tau\), and
\(\ell_u=-\sum_{\tau i}u_{\tau i}X_{\tau i}\).
Suppose every index in the small block \(p\le i<N\) is separated
from every other index of its label by at least \(g>0\).
Let \(F=\ell_u^{d+1}\) and \(D(u)=\deg_{\rm Fl}(\ell_u^d)>0\).
For \(k\ge2,\ m=k-1\), (FA3)--(FA4) give the exact absolute ratio
\[
 \frac{|k\deg_{\rm Fl}(X_{\sigma i}^{m}Q_{\tau k})|}
 {(d+1)D(u)}
 =m!\left|[t^m]\frac{P(u+t e_{\sigma i})}{P(u)}
                    H_{\tau k}(u+t e_{\sigma i})\right|.
 \tag{FA5}
\]
The common sign from \(\ell_u\) has disappeared. For this calculation
alone translate the entries of label \(\tau\) by minus their
minimum. The degree is unchanged and (FA2) kills the difference
of the two linear powers, which is divisible by \(p_{\tau1}\).
Thus (FA5) can be estimated using entries of label \(\tau\)
in \([0,W_\tau]\); this is not a translation of the arithmetic
metric or of its determinant term.

The shifted Vandermonde ratio is a product of \(N-1\) factors
\(1+t/(u_{\sigma i}-u_{\sigma j})\). In absolute coefficients it is
bounded by the exponential with parameter \((N-1)/g\).
The shifted Schur expression is bounded coefficientwise by
\((2^{k-1}/k!)(W_\tau+t)^k\) if \(\sigma=\tau\), and the
constant bound suffices otherwise. Since
\(W_\tau>0\), multiply exponential majorants to obtain
\[
 \frac{|k\deg_{\rm Fl}(X_{\sigma i}^{k-1}Q_{\tau k})|}
 {(d+1)D(u)}
 \le\frac{2^{k-1}}{k!}W_\tau
       \bigl(N+(N-1)W_\tau/g\bigr)^{k-1}.
 \tag{FA6}
\]
Modulo a full power sum, every source \(pLie\) of positive
degree is up to sign the power sum in the small block of one label.
The downward-shift identity justifies this replacement inside
\(\deg_{\rm Fl}\). Thus its bound is \(q\) times (FA6).

For clarity, write \(R_B(F)\) for the original
\texttt{restSum} summed over \(\tau,k\):
\[
 R_B(F)=\sum_{\ell>B,\ {\rm prime},\ \ell\notin{\rm Bad}}
 \log\ell\sum_{\tau,k\ge2}\sum_{a\ge1}
 \ell^{-ka} k\deg_{\rm Fl}
       (pLie(k-1,\mathrm{frob}(\ell,\tau,a))Q_{\tau k}).
 \tag{FA7}
\]
The orbit labels and signs are not changed. The elementary
Chebyshev estimate \(\sum_{\ell\le x}\log\ell\le4\log2\,x<3x\)
and partial summation give, for \(B\ge128,\ k\ge2\),
\[
 \sum_{\ell>B}\frac{\log\ell}{\ell^k-1}
 \le 7B^{1-k}.
 \tag{FA8}
\]
For example the bound for \(\ell^{-k}\) is
\(3kB^{1-k}/(k-1)\), and
\((1-B^{-k})^{-1}\le(1-128^{-2})^{-1}\).
With the actual cutoff condition
\(4fN^2(1+U/g)< B+1\), set
\(x_\tau=2(N+(N-1)W_\tau/g)/B\).
Then \(x_\tau\le129/(128 fN)\le129/1024\).
Summing (FA6)--(FA8) proves the stronger uniform bound
\[
 \boxed{\quad |R_B(\ell_u^{d+1})|\le
   (d+1)D(u)\,qW .\quad}
 \tag{FA9}
\]
Indeed \(7\sum_{m\ge1}x_\tau^m/(m+1)!
\le7x_\tau/(2(1-x_\tau))<1\).
This proves the regular-weight supplier, with the floor convention
on \(B\) left intact.

\paragraph{The singular Grassmannian test.}
Let \(\mu=\sum_{\tau,i\ge p}X_{\tau i}\) and
\[
 V=\prod_\tau\prod_{\substack{i<j\\i,j\ {\rm in\ same\ block}}}
       (X_{\tau j}-X_{\tau i}),\qquad F_0=\mu^{d_0+1}V,\qquad
 J=\deg_{\rm Fl}(\mu^{d_0}V).
\]
Take \(z^0\) equal to \(0\) in the first block and \(1\) in the
second, in each label. Factoring \(P\) into its within-block and
cross-block Vandermondes shows
\(J=(-1)^dd_0! V(\partial)P(z^0)>0\):
only the full within-block degree survives differentiation;
the cross-block sign is \((-1)^{d_0}\), the within-block sign
is \((-1)^{d-d_0}\), and \(V(\partial)V=\sum_\alpha
\alpha!\,([X^\alpha]V)^2>0\).
In particular no lower estimate for \(J\) is required.

Let \(I\) be the small block in the label selected by \(pLie\).
As in (FA5), apply \(V(\partial)\) to the sum of shifted
determinants in (FA4) and divide by \(V(\partial)P(z^0)\).
Define
\[
 R(t,z)=\sum_{i\in I}\frac{P(z+t e_i)}{P(z)}
                         H_{\tau k}(z+t e_i).
\]
The numerator is alternating within each block. Its within-block
Vandermonde factors therefore cancel the corresponding factors
in \(P\); the cross differences are nonzero near \(z^0\).
Thus \(R\) is regular in \(z\) there, polynomial in \(t\).
Only the term where all within-block factors are differentiated
can survive in \(V(\partial)(P R)(z^0)\).
Using (FA3) and \((d_0+1)!/d_0!=d_0+1\) yields
\[
 \frac{|k\deg_{\rm Fl}(pLie(k-1,s)Q_{\tau k})|}
 {(d_0+1)J}=m!\,|[t^m]R(t,z^0)|,\quad m=k-1 .
 \tag{FA10}
\]
Here the sign of the small-block power sum is immaterial.

To bound the confluent divided difference, let
\(P_I(Z)=\prod_{i\in I}(Z-z_i)\) and
\(C(Z,t)=\prod_{j\notin I\ {\rm in\ this\ label}}
(1+t/(Z-z_j))\). If the moved label is \(\tau\),
extend \(H_{\tau k}\) to the virtual alphabet
\(z_\tau+(Z+t)-Z\) by its symmetric-function expression.
If it is not \(\tau\), keep \(H_{\tau k}\) constant.
Call the resulting polynomial \(H^{\rm ext}(Z,t)\).
The residue formula for distinct \(z_i\) is
\[
 R(t,z)=\frac{1}{2\pi\mathrm{i}}\oint
 \frac{P_I(Z+t)}{tP_I(Z)} C(Z,t) H^{\rm ext}(Z,t)\,dZ .
\]
For \(t=0\) take the removable value of the integral.
The contour encloses only the small block. At \(z^0\), write
\(Z=1+\zeta\); taking the residue gives
\[
 R(t,z^0)=\sum_{l=1}^q\binom ql t^{l-1}
 [\zeta^{l-1}]\bigl(C(1+\zeta,t)H^{\rm ext}(1+\zeta,t)\bigr),
 \qquad C(1+\zeta,t)=(1+t/(1+\zeta))^p .
 \tag{FA11}
\]
This evaluates the singular test directly, not by dividing
by its vanishing full Vandermonde.

For \(|\zeta|=1/3,\ |t|\le1/N\), all entries of the virtual
alphabet, including \(Z+t\) and \(Z\), have absolute value at most
\(2\). The usual Schur addition identity for a negative one-letter
alphabet gives
\[
 s_\lambda(X+y-Z)=
 \sum_{\mu\subset\lambda,\ \lambda/\mu\ {\rm vertical\ strip}}
 (-Z)^{|\lambda|-|\mu|}s_\mu(X+y).
\]
This follows from Jacobi--Trudi and Cauchy--Binet using
\(\sum h_i(-Z)T^i=1-ZT\), or from skew Schur addition
and conjugation. The sum of \(s_\mu(1^{N+1})\) in this formula
is at most \(s_\lambda(1^{N+1+k})\) for \(|\lambda|=k\):
extend each tableau of \(\mu\) by filling the omitted cells
of row \(i\) with \(N+1+i\). This is a semistandard tableau,
and the omitted cells and original tableau can be recovered.
The hook-content formula, with cells ordered row by row,
gives
\[
 \frac{s_\lambda(1^{N+1+k})}{s_\lambda(1^N)}
 \le \prod_{j=0}^{k-1}\frac{N+1+k-j}{N-j}
 \le \binom{2k+1}{k}\le 2\cdot4^k ,
\]
since the content of the \(j\)-th cell (starting at zero)
is at least \(-j\) and \(k\le N\).
Thus (FA4) gives
\(|H^{\rm ext}|\le 16^k/k!\), also in the constant case.
By Cauchy's coefficient inequality, the sum in (FA11) is bounded
on \(|t|=1/N\) by
\[
 \frac{16^k}{k!}e^{3p/(2N)}
 q(1+3/N)^{q-1}<13q\,\frac{16^k}{k!}.
\]
We used \(3q+1\le N\).
Another Cauchy estimate, now in \(t\), proves from (FA10)
\[
 \frac{|k\deg_{\rm Fl}(pLie(k-1,s)Q_{\tau k})|}
 {(d_0+1)J}\le 13q\,16^k N^{k-1}/k .
 \tag{FA12}
\]
The source cutoff \(B\ge4fN^2,\ B\ge128\) ensures
\(x=16N/B\le1/2\). Equations (FA7)--(FA8) now give
\[
 \boxed{\quad
 \frac{|R_B(F_0)|}{fq(d_0+1)J}
 \le 7\cdot13\cdot16\sum_{m\ge1}\frac{x^m}{m+1}
 \le728 .\quad}
 \tag{FA13}
\]
This replaces the sharper source constant \(8/(5fq)\) by
a dimension-independent constant sufficient for the weighted-height
bound. No additional geometric assumption is used for either tail.

\paragraph{Effect on the weighted-height comparison.}
Use the actual signed base test
\(\beta_B\ge-107\Lambda_B/2\) from the determinant assembly,
with \(\Lambda_B=\log(B\operatorname{rd}(E))+|\lambda_E|\).
Suppose the same-tower comparison (including its fixed flag
height functional) gives \(h_u\ge0\),
\(h_0\le\frac12\log(N+1)\), and \(|c_B|\le4\Lambda_B\),
in the notation of the tower-intersection reduction.
The algebraic elimination there, now using (FA9) and (FA13), gives
\[
 \sigma_B(u)\ge-236\Lambda_B L(u)-qW .
 \tag{FA14}
\]
Indeed \(Z=(107/2)\Lambda_B+\frac12\log(N+1)+728
\le236\Lambda_B\) because \(\log B>4\), \(N+1\le B\).
Both entries of \(\max(Z,(q/p)Z+(N/p)4\Lambda_B)\)
are bounded by \(236\Lambda_B\).
If the remaining original volume and norm-count identifications hold,
the previously established \(80\Lambda_B L(u)\) allowance gives
the conditional mean bound \(-316\Lambda_B L(u)-qW\).
The tail estimates themselves are unconditional polynomial statements;
the geometric comparison hypotheses in (FA14) remain necessary.
The earlier sharper conditional constant using the sharper test
remainder is not invalidated.
